
%%
%% Lab 6 report content.  The document preamble and package definitions
%% can be retained from lab4.tex.
\title{LAB 6: Axisymmetric and 3D Solid Elements, Time-History Output}

\author{Adhilsha A (SE25PCSE007)}
\maketitle

\tableofcontents

\section{Thick-Walled Cylinder: Problem Definition}

The objective of this experiment is to analyse a long, thick-walled steel
cylinder using an axisymmetric model and a three-dimensional quarter-sector
model.  The pressure is applied through three load steps and the radial
displacement, radial stress, and hoop stress are extracted using the ANSYS
time-history postprocessor.  The finite element results are compared with the
analytical Lam\'e solution.

\begin{table}[h!]
\centering
\begin{tabular}{ll}
\toprule
\textbf{Quantity} & \textbf{Value} \\
\midrule
Inner radius, $r_i$ & $\SI{100}{\milli\metre}$ \\
Outer radius, $r_o$ & $\SI{200}{\milli\metre}$ \\
Slice height, $h$ & $\SI{20}{\milli\metre}$ \\
Young's modulus, $E$ & $\SI{200000}{\mega\pascal}$ \\
Poisson's ratio, $\nu$ & $0.30$ \\
Internal pressure at $t=1$, $p_i$ & $\SI{10}{\mega\pascal}$ \\
External pressure at $t=2$, $p_o$ & $\SI{5}{\mega\pascal}$ \\
Analysis assumption & Linear elastic plane strain \\
\bottomrule
\end{tabular}
\caption{Geometry, material properties, and loading data}
\end{table}

The loading history used in ANSYS was:
\begin{table}[h!]
\centering
\begin{tabular}{ccl}
\toprule
\textbf{Load step} & \textbf{Time} & \textbf{Loading} \\
\midrule
1 & $1$ & $p_i$ ramped from $0$ to $\SI{10}{\mega\pascal}$ \\
2 & $2$ & $p_i=\SI{10}{\mega\pascal}$ held; $p_o$ ramped to $\SI{5}{\mega\pascal}$ \\
3 & $3$ & Both pressures ramped to zero \\
\bottomrule
\end{tabular}
\caption{Three-step pressure history}
\end{table}

Five substeps were used in every load step.  The loads were ramped using
\texttt{KBC,0}, and all substeps were written using
\texttt{OUTRES,ALL,ALL}.

\subsection{Lam\'e Solution}

For a thick cylinder subjected to internal pressure $p_i$ and external
pressure $p_o$, the radial and hoop stresses are
\[
\sigma_r=A-\frac{B}{r^2},\qquad
\sigma_\theta=A+\frac{B}{r^2},
\]
where
\[
A=\frac{p_i r_i^2-p_o r_o^2}{r_o^2-r_i^2},
\qquad
B=\frac{(p_i-p_o)r_i^2r_o^2}{r_o^2-r_i^2}.
\]
For the plane-strain cylinder, the axial stress is $\sigma_z=2\nu A$.
At the bore, the radial and hoop stresses are also obtained directly from
\[
\sigma_r(r_i)=-p_i,\qquad
\sigma_\theta(r_i)=
\frac{p_i(r_o^2+r_i^2)-2p_o r_o^2}{r_o^2-r_i^2}.
\]

For load step 1, $p_i=\SI{10}{\mega\pascal}$ and $p_o=0$, giving
\[
A=\SI{3.3333}{\mega\pascal},\qquad
B=\SI{133333.33}{\mega\pascal\milli\metre^2},
\]
\[
\boxed{
\sigma_r(r_i)=-\SI{10}{\mega\pascal},\qquad
\sigma_\theta(r_i)=\SI{16.6667}{\mega\pascal}.}
\]

For load step 2, $p_i=\SI{10}{\mega\pascal}$ and
$p_o=\SI{5}{\mega\pascal}$, giving
\[
A=-\SI{3.3333}{\mega\pascal},\qquad B=\SI{66666.67}{\mega\pascal\milli\metre^2},
\]
\[
\boxed{
\sigma_r(r_i)=-\SI{10}{\mega\pascal},\qquad
\sigma_\theta(r_i)=\SI{3.3333}{\mega\pascal}.}
\]

At load step 1, at the outer radius, $(r = r_o)$, \text{ the hoop stress is given by:}
\[
\sigma_\theta(r_o) = A + \frac{B}{r_o^2} = \frac{2 p_i r_i^2 - p_o (r_o^2 + r_i^2)}{r_o^2 - r_i^2}
\]
\[
\boxed{\sigma_\theta(r_o) = 6.6667\,\text{MPa}}
\]

Using $E = 200\,000\text{ MPa}$ and $\nu = 0.30$, the analytical plane-strain radial displacement is given by
\begin{equation}
u_r(r) = \frac{1 + \nu}{E} \left[ (1 - 2\nu)Ar + \frac{B}{r} \right].
\end{equation}

\textbf{At $t = 1$ ($A = 3.3333\text{ MPa}$, $B = 133\,333.33\text{ MPa}\cdot\text{mm}^2$):}
    \begin{align*}
    u_r(r_i) &= \frac{1.3}{200\,000} \left[ 0.4(3.3333)(100) + \frac{133\,333.33}{100} \right] \\
             &= 6.5 \times 10^{-6} \times [133.33 + 1333.33] 
    \end{align*}
    \[\boxed{u_r(r_i)  = {0.009533\text{ mm}}}
    \]

\textbf{At $t = 2$ ($A = -3.3333\text{ MPa}$, $B = 66\,666.67\text{ MPa}\cdot\text{mm}^2$):}
    \begin{align*}
    u_r(r_i) &= \frac{1.3}{200\,000} \left[ 0.4(-3.3333)(100) + \frac{66\,666.67}{100} \right] \\
             &= 6.5 \times 10^{-6} \times [-133.33 + 666.67] 
    \end{align*}
    \[\boxed{u_r(r_i)  = {0.003467\text{ mm}}}\]

\section{Axisymmetric Model}

\subsection{Model and APDL Code}

The meridional $r$--$z$ slice was modelled using the eight-node
\texttt{PLANE183} element in axisymmetric mode.  The horizontal coordinate
represents radius and the vertical coordinate represents the cylinder axis.
The top and bottom edges were constrained in the axial direction
($U_Y=0$), which represents the plane-strain condition.

\lstinputlisting[style=pyconsole]{codes/lab6_1.txt}

\subsection{Mesh and Deformed Configuration}

The mapped quadrilateral mesh of the axisymmetric slice and the deformed configuration is shown in
Figure~\ref{lab6_1_mesh}.  The expanded axisymmetric representation is shown
in Figure~\ref{lab6_1_axis} for visualization.

\subsection{Stress Contours}

At $t=1$, the radial stress varies from approximately
$-\SI{10}{\mega\pascal}$ at the bore to nearly zero at the outer surface.
The hoop stress is maximum at the bore, with a value close to
$\SI{16.67}{\mega\pascal}$, and decreases to $\SI{6.67}{\mega\pascal}$ towards the outer radius.
The corresponding radial and hoop stress contours are shown in
Figures~\ref{lab6_1_radial} and \ref{lab6_1_hoop}.

\subsection{Hoop Stress Along the Radius}

The hoop stress was mapped onto a path from $r=100$ mm to $r=200$ mm and is shown in Figure~\ref{lab6_1_path}.  The stress decreases
monotonically from approximately $\SI{16.67}{\mega\pascal}$ at the bore to
$\SI{6.67}{\mega\pascal}$ at the outer surface, in agreement with the Lam\'e
solution.

\subsection{POST26 Time Histories}

The bore node was selected using \texttt{NB=NODE(100,0,0)}.  The radial
displacement was obtained using \texttt{NSOL}, while the radial and hoop
stresses were obtained using \texttt{ESOL}.  The plotted histories are shown
in Figures~\ref{lab6_1_history_stress} and \ref{lab6_1_history_disp}.

The POST26 listing at different time steps is given in
Figure~\ref{lab6_1_listing}.  The principal values read from the histories are
summarized below in Table~\ref{tab1}.  Small differences from the analytical values (e.g. 10 MPa) are due to
finite element interpolation and the selected element/node location.

% \begin{table}[h!]
% \caption{Axisymmetric bore values from POST26 and the analytical solution}
% \label{tab1}
% \centering
% \begin{tabular}{c|ccc|ccc}
% \toprule
% & \multicolumn{3}{c|}{\textbf{Axisymmetric FE}} 
% % & \multicolumn{3}{c}{\textbf{Lam\'e / reference}}\\
% \textbf{Time} & $u_r$ (mm) & $\sigma_r$ (MPa) & $\sigma_\theta$ (MPa)
% % & $u_r$ (mm) & $\sigma_r$ (MPa) & $\sigma_\theta$ (MPa)\\
% \midrule
% 1 & $0.0095$ & $-9.9998$ & $16.6667$
%   % & $0.0087$ & $-10.00$ & $16.67$\\
% 2 & $0.0035$ & $-9.9998$ & $6.6667$
%   % & $0.0043$ & $-10.00$ & $6.67$\\
% 3 & $0.0000$ & $0.0000$ & $0.0000$
%   % & $0.0000$ & $0.00$ & $0.00$\\
% \bottomrule
% \end{tabular}
% \end{table}

\begin{table}[h!]
\caption{Axisymmetric bore values from POST26}
\label{tab1}
\centering
\begin{tabular}{c|ccc}
\toprule
& \multicolumn{3}{c}{\textbf{Axisymmetric FE}} \\ % & \multicolumn{3}{c}{\textbf{Lam\'e / reference}}\\
\textbf{Time} & $u_r$ (mm) & $\sigma_r$ (MPa) & $\sigma_\theta$ (MPa) \\ % & $u_r$ (mm) & $\sigma_r$ (MPa) & $\sigma_\theta$ (MPa)\\
\midrule
1 & $0.0095$ & $-9.9998$ & $16.6667$ \\ % & $0.0087$ & $-10.00$ & $16.67$\\
2 & $0.0035$ & $-9.9999$ & $3.3332$ \\ % & $0.0043$ & $-10.00$ & $6.67$\\
3 & $0.0000$ & $0.0000$ & $0.0000$ \\ % & $0.0000$ & $0.00$ & $0.00$\\
\bottomrule
\end{tabular}
\end{table}



\subsection{Hoop Stress Versus Internal Pressure}

% For load step 1, the external pressure is zero and
% \[
% \sigma_\theta(r_i)
% =\frac{r_o^2+r_i^2}{r_o^2-r_i^2}p_i
% =\frac{200^2+100^2}{200^2-100^2}p_i
% =1.6667p_i.
% \]
% Thus, the expected slope of the hoop-stress-versus-internal-pressure plot is
% approximately $1.667$.  
The POST26 plot in
Figure~\ref{lab6_1_pressure} is linear and The slope $m$ is calculated using two points $(2,\, 3.3333)$ and $(10,\, 16.6667)$:
\[
m = \frac{\Delta \sigma_\theta}{\Delta p_i} = \frac{16.6667 - 3.3333}{10 - 2} = \frac{13.3333}{8} \approx 1.6667
\]
This matches the theoretical stress concentration factor at the inner radius$(r = r_i)$ for $p_o = 0$:
\[
m = \frac{r_o^2 + r_i^2}{r_o^2 - r_i^2} = \frac{200^2 + 100^2}{200^2 - 100^2}  \approx 1.6667
\]




During load step 2, the internal pressure remains fixed, so the radial stress
at the bore remains $\sigma_r=-p_i=-\SI{10}{\mega\pascal}$.  The external
pressure changes the stress constant $A$ and reduces the pressure difference
across the wall.  Consequently, the hoop stress at the bore decreases from
approximately $\SI{16.67}{\mega\pascal}$ to $\SI{6.67}{\mega\pascal}$.

\begin{figure}[!h]
    \centering
    \includegraphics[width=0.78\linewidth]{figures/lab6_1_0.png}
    \caption{Mapped mesh of the axisymmetric rectangular slice}
    \label{lab6_1_mesh}
\end{figure}

\begin{figure}[!h]
    \centering
    \includegraphics[width=0.78\linewidth]{figures/lab6_1_0_axis.png}
    \caption{Expanded axisymmetric representation of the cylinder}
    \label{lab6_1_axis}
\end{figure}
\begin{figure}[!h]
    \centering
    \includegraphics[width=0.73\linewidth]{figures/lab6_1_1_radial.png}
    \caption{Axisymmetric radial stress, $\sigma_r$, at $t=1$}
    \label{lab6_1_radial}
\end{figure}

\begin{figure}[!h]
    \centering
    \includegraphics[width=0.73\linewidth]{figures/lab6_1_2_hoop.png}
    \caption{Axisymmetric hoop stress, $\sigma_\theta$, at $t=1$}
    \label{lab6_1_hoop}
\end{figure}

\begin{figure}[!h]
    \centering
    \includegraphics[width=0.73\linewidth]{figures/lab6_1_3_hoop_along_radius.png}
    \caption{Hoop stress along the radial path at $t=1$}
    \label{lab6_1_path}
\end{figure}

\begin{figure}[!h]
    \centering
    \includegraphics[width=0.72\linewidth]{figures/lab6_1_4_time2_var23.png}
    \caption{Bore radial and hoop stresses against time}
    \label{lab6_1_history_stress}
\end{figure}

\begin{figure}[!h]
    \centering
    \includegraphics[width=0.72\linewidth]{figures/lab6_1_4_time2_var2.png}
    \caption{Bore radial displacement against time}
    \label{lab6_1_history_disp}
\end{figure}
\begin{figure}[!h]
    \centering
    \includegraphics[width=0.73\linewidth]{figures/lab6_1_4_time2_prvar234.png}
    \caption{POST26 listing of bore displacement and stresses}
    \label{lab6_1_listing}
\end{figure}
\begin{figure}[h]
    \centering
    \includegraphics[width=0.72\linewidth]{figures/lab6_1_5_last2.png}
    \caption{Bore hoop stress versus internal pressure for load step 1}
    \label{lab6_1_pressure}
\end{figure}


\section{Three-Dimensional Quarter-Sector Model}

\subsection{Model and APDL Code}

The same cylinder slice was represented by a 90-degree sector using
\texttt{SOLID186} twenty-node brick elements.  Symmetry boundary conditions
were applied on the two cut faces, and $U_Z=0$ was imposed on the top and
bottom faces.  The pressure was applied to the inner and outer cylindrical
faces in the cylindrical coordinate system.

\lstinputlisting[style=pyconsole]{codes/lab6_2.txt}

\subsection{Brick Mesh and Deformed Shape}

The swept brick mesh and the deformed shape is shown in Figure~\ref{lab6_2_brick}.  The and the hoop-stress contour at $t=1$ are shown in
Figure~\ref{lab6_2_hoop}.  The stress field is circumferentially symmetric,
as expected for uniform pressure and the two symmetry boundaries. The values are similar to the ones from before, but with slightly lesser accuracy. The mesh had 8000 elements, 24481 nodes, and 48962 DOF.

\begin{figure}[!h]
    \centering
    \includegraphics[width=0.72\linewidth]{figures/lab6_2_0_brick.png}
    \caption{Swept brick mesh of the 90-degree cylinder sector}
    \label{lab6_2_brick}
\end{figure}

\begin{figure}[!h]
    \centering
    \includegraphics[width=0.72\linewidth]{figures/lab6_2_1_hoop.png}
    \caption{Three-dimensional hoop stress contour at $t=1$}
    \label{lab6_2_hoop}
\end{figure}

The model contains three translational degrees of freedom per node.  Hence,
\[
\mathrm{DOF}_{3D}=3\,N_{\mathrm{nodes}}.
\]
The exact values from the simulation gave  us 320 elements, 1973 nodes, and 5919 total degrees of freedom (DOF).
This is larger than the axisymmetric model because the circumferential
direction is discretized explicitly.

\subsection{POST26 Histories}

The bore node on the global $x$-axis was selected using
\texttt{NB=NODE(100,0,0)}.  This location is convenient because its radial
direction coincides with global $x$ and its hoop direction coincides with
global $y$.  Therefore, in the global stress output,
\[
\sigma_r=S_X,\qquad \sigma_\theta=S_Y
\]
at this node.  The resulting histories are shown in
Figure~\ref{lab6_2_history} and Figure~\ref{lab6_2_listing}. These matches closely with the results from before.

\begin{figure}[!h]
    \centering
    \includegraphics[width=0.72\linewidth]{figures/lab6_2_2_radial-n-hoop-time.png}
    \caption{Radial and hoop stresses against time for solid quarter model}
    \label{lab6_2_history}
\end{figure}
\begin{figure}[!h]
    \centering
    \includegraphics[width=0.73\linewidth]{figures/lab6_2_3_table.png}
    \caption{POST26 listing of bore displacement and stresses for solid quarter model}
    \label{lab6_2_listing}
\end{figure}
\subsection{Comparison of Axisymmetric, 3D, and Lam\'e Results}

\begin{table}[h!]
\centering
\begin{tabular}{c|c|ccc|ccc}
\toprule
& & \multicolumn{3}{c|}{$t=1$} & \multicolumn{3}{c}{$t=2$}\\
\textbf{Model} & \textbf{DOF}
& $\sigma_r$ & $\sigma_\theta$ & $u_r$
& $\sigma_r$ & $\sigma_\theta$ & $u_r$\\
& & \multicolumn{2}{c|}{MPa} & mm
& \multicolumn{2}{c|}{MPa} & mm\\
\midrule
Axisymmetric PLANE183 & $48,962$
& $-9.9998$ & $16.6665$ & $0.009533$
& $-9.9999$ & $3.3332$ & $0.003467$\\
3D SOLID186 & $5,919$
& $-9.9648$ & $16.6315$ & $0.009533$
& $-9.9824$ & $3.3157$ & $0.003467$\\
Lam\'e solution & -- 
& $-10.00$ & $16.67$ & $0.009533$
& $-10.00$ & $6.67$ & $0.003467$\\
\bottomrule
\end{tabular}
\caption{Bore response comparison; stresses are in MPa and displacement is in mm}
\end{table}

The three-dimensional results agree closely with both the axisymmetric
solution and the Lam\'e solution.  The small difference in the radial stress
and hoop stress is attributable to mesh interpolation and the use of a
finite sector. With a very high number of elements used in modeling the axisymmetric model, we got a very refined result much closer to the analytical value as seen in the above table.

An axisymmetric model is sufficient when the geometry,
material, loading, and boundary conditions are all rotationally symmetric.
A full three-dimensional model is required when the cylinder has local
features, non-axisymmetric loads, circumferential variation, non-symmetric
supports, holes, contacts, or other three-dimensional effects. 

\section{Optional Tetrahedral Mesh}

For the optional comparison, the swept brick mesh was cleared and the sector
was remeshed using free \texttt{SOLID187} tetrahedral elements.  The resulting
mesh is shown in Figure~\ref{lab6_2_tet}.  The hoop stress contour and the
corresponding time-history plot are shown in
Figures~\ref{lab6_2_tet_hoop} and \ref{lab6_2_tet_history} which match the histories from the last question as well.

\begin{figure}[!h]
    \centering
    \includegraphics[width=0.72\linewidth]{figures/lab6_2_4_0_tetrahedral_mesh.png}
    \caption{Free tetrahedral mesh using SOLID187}
    \label{lab6_2_tet}
\end{figure}

\begin{figure}[!h]
    \centering
    \includegraphics[width=0.72\linewidth]{figures/lab6_2_4_1_hoop.png}
    \caption{Hoop stress contour for the tetrahedral mesh at $t=1$}
    \label{lab6_2_tet_hoop}
\end{figure}

\begin{figure}[!h]
    \centering
    \includegraphics[width=0.72\linewidth]{figures/lab6_2_4_2_radial-n-hoop-time.png}
    \caption{Tetrahedral-model radial and hoop stress histories}
    \label{lab6_2_tet_history}
\end{figure}

The tetrahedral results are close to the brick results, although the free
mesh does not preserve the structured radial and circumferential alignment of
the swept mesh.  Consequently, a finer tetrahedral mesh is generally needed
to obtain the same accuracy near the bore. The listing used for the
tetrahedral comparison is shown in
Figure~\ref{lab6_2_tet_listing}.

\begin{figure}[!h]
    \centering
    \includegraphics[width=0.72\linewidth]{figures/lab6_2_4_3_table.png}
    \caption{POST26 listing for the tetrahedral model}
    \label{lab6_2_tet_listing}
\end{figure}

\section{Discussion and Conclusion}

The axisymmetric and three-dimensional analyses produced the expected
thick-cylinder response.  The radial stress at the bore is controlled by the
internal pressure, while the hoop stress is highest at the bore and decreases
with radius.  When the external pressure is introduced in load step 2, the
radial stress at the bore remains fixed because the internal pressure is held
constant, whereas the hoop stress decreases because the pressure difference
across the wall is reduced.

The numerical hoop-stress distribution follows the Lam\'e solution closely,
and the POST26 pressure plot confirms the theoretical slope
\[
\frac{r_o^2+r_i^2}{r_o^2-r_i^2}=1.6667.
\]
The 3D brick model reproduces the axisymmetric results at the bore when the
node on the global $x$-axis is used, since the global stress components then
correspond directly to the radial and hoop directions.  The tetrahedral model
also gives the same qualitative and quantitative response, but the structured
brick mesh is more efficient for this regular sector geometry.  Overall, the
experiment demonstrates the equivalence of axisymmetric and suitably
constrained 3D models for a rotationally symmetric pressure vessel and shows
how ANSYS POST26 can be used to examine the complete load history.